\section{Data}
\label{sec:data}

\subsection{Preference pairs}
HPDv2 training pairs whose prompt text matches ``tattoo'' contribute \nHPDpairsAll\ pairs, \nHPDpromptsAll\ prompts and \nHPDimgs\ images.
ImageRewardDB training and validation rows with a tattoo prompt contribute \nIRrows\ rows.
\nIRempty\ of those image files are empty in the upstream archive, leaving \nIRimgs\ usable images.
Pick-a-Pic v2 rows from a public parquet mirror contribute \nPAProws\ tattoo rows, of which \nPAPpairs\ are labelled and show two different images, across \nPAPprompts\ prompts.
The original Hugging Face dataset card has been removed.
Mirrors state an MIT licence.
We could not verify that statement, so the mirror is used only as local metadata and local images, and we do not redistribute either.
Keyword selection keeps \nPAPkwPairs\ labelled pairs over \nPAPkwPrompts\ prompts.
Images for those rows, \nPAPimgs\ files, were read by row-group ranges from the mirror.
CLIP then keeps only design and on-skin pairs.

\subsection{Tattoo-subject filter}
Keyword rules mark an explicit design request or an incidental ``person with tattoos'' prompt.
Otherwise, and when the keyword and the image disagree, we use CLIP ViT-B/32 zero-shot among four captions: a design on paper, a close-up of a tattoo on skin, a picture of a person who has tattoos, and something else \citep{radford2021clip}.
The automatic labels are design, on-skin, incidental, other, or ambiguous.
Counts before review are in Table~\ref{tab:data}.

Ambiguous files, \nAmb\ images in total, were then labelled by a single non-blinded rater, the author's assistant, who saw the prompt and the image.
The assistant's labels were design \NambDesign, on-skin \NambOnskin, incidental \NambIncidental\ and other \NambOther.
They are not a second independent annotation, and they were not blinded to the study aim.
After that replacement, tattoo-subject images are those labelled design or on-skin.
Preference probes and all four reward models on the reviewed HPDv2 tattoo-subject subset use the same \HpdHpstwoN\ pairs over \NHpdTsprompts\ prompts.
ImageRewardDB uses \IrFractightgapfiveN\ pairs over \NIrTsprompts\ prompts.
The Pick-a-Pic tattoo-subject slice that remains after the CLIP filter is \PapHpstwoN\ pairs over \NPapTsprompts\ prompts.
That slice is very small.
ImageRewardDB tattoo-subject has only \NIrTsprompts\ prompts.
Its intervals are wide because the bootstrap resamples prompts, not pairs.

\subsection{Defect corpus}
Controlled defects and the spread test use \nDegrDesigns\ designs drawn with seed 0 from Drozdik/tattoo\_v0.
That set is released for research under CC BY-NC-SA.
We do not redistribute those images.
Six operators at two levels produce \nDegrImgs\ images including the clean copies.
Reward scoring of defects uses \nDegrRewDesigns\ of the designs.

% auto-generated by tables/make_tables.py
\begin{table}[t]
\centering
\caption{Corpora in this audit. Tattoo-subject images are design or on-skin. HPDv2 and ImageRewardDB ambiguous files (\nAmb\ images) were relabelled by one non-blinded rater. Pick-a-Pic v2 images were fetched and then CLIP-filtered; only design and on-skin pairs were scored. Pair counts are in Table~\ref{tab:agree}.}
\label{tab:data}
\footnotesize
\resizebox{\linewidth}{!}{%
\begin{tabular}{lrrr}
\toprule
 & HPDv2 & ImageRewardDB & Pick-a-Pic v2 \\
\midrule
Prompts with ``tattoo'' & \nHPDpromptsAll & --- & \nPAPprompts \\
Images or rows & \nHPDimgs & \nIRimgs\ / \nIRrows & \nPAProws \\
Images fetched & --- & --- & \nPAPimgs \\
Empty files dropped & --- & \nIRempty & --- \\
Design, before review & \NHpdClsDesign & \NIrClsDesign & \NPapClsDesign \\
Ambiguous, before review & \NHpdClsAmbiguous & \NIrClsAmbiguous & \NPapClsAmbiguous \\
Design, after review & \NHpdClstwoDesign & \NIrClstwoDesign & --- \\
On skin, after review & \NHpdClstwoOnskin & \NIrClstwoOnskin & --- \\
Incidental, after review & \NHpdClstwoIncidental & \NIrClstwoIncidental & --- \\
Other, after review & \NHpdClstwoOther & \NIrClstwoOther & --- \\
Keyword pairs / prompts & --- & --- & \nPAPkwPairs\ / \nPAPkwPrompts \\
Tattoo-subject prompts & \NHpdTsprompts & \NIrTsprompts & \NPapTsprompts \\
\bottomrule
\end{tabular}}
\end{table}

